% !TeX program = xelatex
\documentclass[11pt,a4paper]{article}
\usepackage{fontspec}
\newfontfamily\greekfont[Script=Greek]{Walleye}
\newfontfamily\greekfontsf[Script=Greek]{Walleye}
\usepackage{polyglossia}
\setdefaultlanguage{greek}
\usepackage{enumitem}
\usepackage{amssymb,amsmath,amsthm,booktabs,changepage,diagbox,fontawesome,subcaption,tikz,xcolor,xfrac}
\usepackage[margin=1in]{geometry}
\usepackage{mathpazo}
\usepackage[colorlinks=true, urlcolor=blue]{hyperref}

\usetikzlibrary{calc}
\usetikzlibrary{intersections}
\usetikzlibrary{patterns}
\usetikzlibrary{decorations.footprints}

\pgfdeclarelayer{bg}
\pgfsetlayers{bg,main}

\renewcommand{\thesubfigure}{\roman{subfigure}}

\newcommand{\abs}[1]{{\left|#1\right|}}
\newcommand{\pn}[1]{{\left(#1\right)}}
\newcommand{\qn}[1]{{\left[#1\right]}}
\newcommand{\mn}[1]{{\left\{#1\right\}}}
\newcommand{\R}{\mathbb{R}}
\newcommand{\Q}{\mathbb{Q}}
\newcommand{\Z}{\mathbb{Z}}
\newcommand{\N}{\mathbb{N}}
\newcommand{\tor}{{\text{ ή }}}
\newcommand{\tand}{{\text{ και }}}
\newcommand{\tif}{{\text{ αν }}}

\definecolor{twitter}{RGB}{29,161,242}
\definecolor{wordpress}{RGB}{71,167,200}
\definecolor{instagram}{RGB}{240,76,92}
\definecolor{facebook}{RGB}{66,103,178}

\let\sin\relax
\DeclareMathOperator{\sin}{\text{{ημ}}}

\let\cos\relax
\DeclareMathOperator{\cos}{\text{{συν}}}

\let\tan\relax
\DeclareMathOperator{\tan}{\text{{εφ}}}

\let\cot\relax
\DeclareMathOperator{\cot}{\text{{σφ}}}

\DeclareMathOperator{\sgn}{sgn}

\theoremstyle{definition}

\newtheorem{them}{Θέμα}
\newtheorem{prop}{Πρόταση}[section]
\newtheorem{cor}{Πόρισμα}[section]
\newtheorem{quest}{Απορία}[section]

\everymath{\displaystyle}

\title{\Huge{\bf Catch Me If You Can}\\{\Large'Η πώς να σκοτώσετε την ώρα σας στην καραντίνα}}
\author{{\color{facebook}\faFacebook}\ \href{www.facebook.com/aftermaths.gr}{/aftermaths.gr} 
{\color{instagram}\faInstagram}\ \href{https://www.instagram.com/aftermathsgr}{@aftermathsgr} {\color{twitter}\faTwitter}\,\href{https://twitter.com/aftermathsgr}{@aftermathsgr}}
\date{\today}

\begin{document}
	\maketitle
	\section{Τρέχει, τρέχει, τρέχει το νερό\ldots}
	Πέρα όμως απ' το νερό, τρέχουν κι άλλα πράγματα: ο χρόνος, οι υποχρεώσεις, η ζωή, τα προβλήματα, οι καθημερινότητα, τα παιδιά στο πάρκο, οι σκέψεις μας\ldots Αλλά, εμάς δε μας απασχολούν όλα αυτά. Εμάς μας ενδιαφέρουν τα μαθηματικά και, για την ακρίβεια, αυτήν τη φορά μας ενδιαφέρουν οι ταχύτητες. Θέτουμε ευθύς αμέσως ένα απλό, άμεσο, πλην όμως μπαμπέσικο, ερώτημα:
	\begin{quest}\label{quest:velocity}
		Γίνεται να τρέχω συστηματικά με μεγαλύτερη ταχύτητα από τον μπροστινό μου και να μην καταφέρω να τον περάσω ποτέ;
	\end{quest}
	Χαζή ερώτηση, θα έλεγε κανείς με μία πρώτη ανάγνωση. Αν πηγαίνεις πιο γρήγορα από τον μπροστινό σου, ε, του κερατά, κάποτε θα τον προσπεράσεις. Άλλωστε, πηγαίνεις πιο γρήγορα (ε;).
	
	Για μια στιγμή, όμως, θα πει κανείς εδώ. Αν είμαι 30 μέτρα πίσω από τον μπροστινό μου σε έναν εξαντλητικό μαραθώνιο και ο μπροστινός μου είναι ένα βήμα πριν τον τερματισμό, τότε, εκτός συγκλονιστικού απροόπτου, μάλλον δεν πρόκειται να τον περάσω. Ή, για να το δούμε και λίγο πιο φορμαλιστικά, αν $f(t)$ είναι η συνάρτηση θέσης μου και $g(t)$ είναι η συνάρτηση θέσης του μπροστινού μου, τότε η υπόθεση ότι βρίσκομαι πίσω του μεταφράζεται στα μαθηματικά ως εξής:
	\begin{equation}\label{eq:behind}
	f(t)<g(t).
	\end{equation}
	Επίσης, θεωρώντας ότι οι συναρτήσεις θέσης και των δυό μας είναι αρκετά ομαλές και, για την ακρίβεια, παραγωγίσιμες, η υπόθεση ότι η ταχύτητά μου είναι σταθερά μεγαλύτερη από τη δική του μεταφράζεται ως εξής:
	\begin{equation}\label{eq:ahead}
	f'(t)>g'(t).
	\end{equation}
	Τώρα, το καυτό ζήτημα στο παραπάνω είναι, κατά πώς φαίνεται πού «ζει» το $t$. Διότι, αν για παράδειγμα $t\in(0,\sfrac{1}{2})$ και $f(t)=t^2$ ενώ $g(t)=t$ τότε, ενώ ικανοποιούνται οι Εξισώσεις \eqref{eq:behind} και \eqref{eq:ahead}, δεν έχουμε καμία ελπίδα να προσπεράσουμε ποτέ τον μπροστινό μας, καθώς ο «αγώνας» τελειώνει πριν προλάβουμε να κεφαλαιοποιήσουμε όλη αυτή την παραπανίσια ταχύτητά μας.
	
	Τι γίνεται, όμως, αν δώσουμε στους δρομείς μας άπειρο χρόνο για να κινηθούν; Αν, δηλαδή, θεωρήσουμε δύο παραγωγίσιμες συναρτήσεις $f,g$ που ικανοποιούν τις \eqref{eq:behind} και \eqref{eq:ahead} και είναι ορισμένες στο διάστημα $(0,+\infty)$ τότε μήπως έχει αρκετό χρόνο ο πίσω δρομέας για να περάσει τον μπροστινό του; Η αλήθεια είναι πως, αν κανείς ακούσει ότι ο πίσω δρομέας τρέχει με μεγαλύτερη ταχύτητα από τον μπροστινό του και έχουν στη διάθεσή τους άπειρο χρόνο για να τρέξουν (θεωρώντας ότι η πείνα και η ζωή η ίδια δεν τους φθείρει καθόλου), ε, μάλλον κάποια στιγμή θα καταφέρει ο γρηγορότερος δρομέας να προσπεράσει τον συναγωνιστή του.
	
	Είναι όμως έτσι;
	
	\section{Λίγη σκέψη δε βλάπτει ποτέ}
	Θέλουμε να αποδείξουμε, αν είναι εφτικτό, φυσικά, ότι αν έχουμε δύο συναρτήσεις που ικανοποιούν τις Εξισώσεις \eqref{eq:behind} και \eqref{eq:ahead} και είναι αμφότερες ορισμένες στο $(0,+\infty)$ τότε υπάρχει κάποιο $t_0>0$ έτσι ώστε για $t>t_0$ να ισχύει ότι $f(t)>g(t)$. Με άλλα λόγια, ότι αν ένα δρομέας τρέχει πιο γρήγορα από τον μπροστινό του δρομέα και έχει άπειρο χρόνο στη διάθεσή του, κάποια στιγμή θα τον προσπεράσει.
	
	Αρχικά, ας κάνουμε ένα κόλπο που θα κάνει τη ζωή μας λίγο πιο εύκολη και το πρόβλημά μας ακόμα πιο απλό στη διατύπωσή του. Θεωρούμε τη συνάρτηση $h(t)=f(t)-g(t)$, οπότε, οι παραπάνω σχέσεις γράφονται ως εξής:
	\begin{equation}\label{eq:h_behind}
	h(t)<0,
	\end{equation}
	και:
	\begin{equation}\label{eq:h_ahead}
	h'(t)>0.
	\end{equation}
	Επίσης, το ζητούμενο τώρα είναι να δείξουμε ότι υπάρχει ένα $t_0\in(0,+\infty)$ έτσι ώστε για κάθε $t>t_0$ να ισχύει ότι $h(t)>0$. Αν ξεχάσουμε λίγο την αναλογία με τους δύο δρομείς και υποθέσουμε ότι έχουμε έναν δρομέα του οποίου η θέση περιγράφεται από την $h$, τότε το παραπάνω μας λέει το εξής: «Αν ένας δρομέας βρίσκεται αριστερά από την θέση $O$ και κινείται με θετική ταχύτητα (προς αυτή τη θέση) τότε κάποια στιγμή θα την περάσει».
	
	Ακόμα πιο λογικό, ε; Αν προχωράμε προς τα κάπου και κινούμαστε συνεχώς προς τα εκεί για όσο θέλουμε, τότε θα φτάσουμε μέχρι εκεί. Πόσο πιο προφανές; Κι όμως, ας σκεφτούμε μία συνάρτηση με γραφική παράσταση σαν αυτή που φαίνεται στο Σχήμα \ref{fig:1}.
	
	\begin{figure}[!htb]
		\centering
		\begin{tikzpicture}
		\draw[thick, ->] (-1,0) -- (6,0);
		\draw[thick, ->] (0,-4) -- (0,2);
		\draw[thick, domain={0.25}:{5.5}, samples=100] plot (\x,{-1/\x});
		\end{tikzpicture}
		\caption{Μία καθημερινή συνάρτηση\ldots}
		\label{fig:1}
	\end{figure}

	Αυτή δεν είναι άλλη παρά η γνωστή και μη εξαιρετέα $h(t)=-\sfrac{1}{t}$. Αν την παρατηρήσουμε καλά, είναι ορισμένη στο $(0,+\infty)$, είναι αρνητική και είναι και γνησίως αύξουσα, δηλαδή $h'(t)>0$. Συνεπώς, αν η θέση του δρομέα μας περιγράφεται από την $h$ τότε αυτός θα τρέχει, θα τρέχει, θα τρέχει, πλησιάζοντας ολοένα και περισσότερο τη θέση που τον ενδιαφέρει και δε θα τη φτάνει ποτέ! Και για να μην έχουμε παρεξηγήσεις και να μη δίνουμε και στο δρομέα μας φρούδες ελπίδες, θα μπορούσαμε κάλλιστα να θεωρήσουμε ότι η θέση του συναρτήσει του χρόνου δίνεται από την $h(t)=-1-\sfrac{1}{t}$, οπότε και θα είχαμε τον δρομέα μας να πλησιάζει, αλλά όχι ασυμπτωτικά τη θέση των ονείρων του.
	
	Δηλαδή, δεν μπορούμε να αποδείξουμε αυτό που, με μια πρώτη ματιά, φαίνεται εύλογο, απλά και μόνο διότι δεν ισχύει. Μπορούμε να τρέχουμε διαρκώς προς τον στόχο μας χωρίς ωστόσο να τον φτάνουμε ποτέ. Αν καθίσουμε λίγο και το σκεφτούμε αυτό δεν είναι παράλογο. Με δεδομένο αυτό, μπορούμε αρκετά φυσιολογικά να συμπεράνουμε ότι \textit{επιβραδύνουμε}. Πράγματι, αν παρατηρήσετε την παραπάνω συνάρτηση, θα δείτε ότι είναι \textit{κοίλη}, κάτι που εύκολα επιβεβαιώνεται και με λίγες πράξεις, αφού:
	\[h''(t)=\pn{\frac{1}{t^2}}'=-\frac{2}{t^3}<0.\]
	
	Δηλαδή, η επιτάχυνση του δρομέα μας (η δεύτερη παράγωγος της συνάρτησης της θέσης του) είναι αρνητική κάτι που αντιστοιχεί σε μία παρατηρούμενη επιβράδυνση. Και για να γυρίσουμε λίγο και στο αρχικό μας πρόβλημα, στις Εξισώσεις, δηλαδή, \eqref{eq:behind} και \eqref{eq:ahead}, αν θεωρήσουμε ότι $g(t)=t$, τότε η $f(t)=t-\sfrac{1}{t}$ είναι μία συνάρτηση που έχει όσες ιδιότητες θέλουμε:
	\begin{enumerate}
		\item Ικανοποιεί τη σχέση $f(t)<g(t)$ για κάθε $t>0$.
		\item Ικανοποιεί τη σχέση $f'(t)>g'(t)$ για κάθε $t>0$ αφού $f'(t)=1+\sfrac{1}{t^2}$ και $g'(t)=1$.
	\end{enumerate}
	Συνεπώς, αν διαλέξουμε μία συνάρτηση σαν την $h$, δηλαδή αρνητική και γνησίως αύξουσα και μία οποιαδήποτε συνάρτηση $g$ θέλουμε, μπορούμε να φτιάξουμε ένα ζευγάρι συναρτήσεων που να ικανοποιεί τις παραπάνω «παράδοξες» σχέσεις.
	
	\section{Γενικά;}
	Αφήνοντας για λίγο κατά μέρους τις φυσικές ερμηνείες που δώσαμε στο παραπάνω πρόβλημα, θα κάνουμε μία απόπειρα να γενικεύσουμε αυτό το ενδιαφέρον και «παραδοξολογικό» αποτέλεσμα. Με σχετική ευκολία μπορούμε να βρούμε συναρτήσεις οι οποίες να γεμίζουν και τα τέσσερα κελιά του Πίνακα \ref{tab:1}.
	
	\begin{table}[!htb]
		\centering
		\begin{tabular}{c|c|c}
			\diagbox{$f$}{$f'$} & $+$ & $-$ \\\hline
			$+$ & \textbf{??} & \textbf{??} \\\hline
			$-$ & \textbf{??} & \textbf{??} 
		\end{tabular}
		\caption{Συνδυασμοί προσήμων.}
		\label{tab:1}
	\end{table}

	Με άλλα λόγια, το πρόσημο μίας συνάρτησης δε μας λέει, άμεσα, τίποτα για το πρόσημο της παραγώγου της και αντίστροφα. Ως εκ τούτου, μπορούμε να βρούμε αρκετά εύκολα θετικές συναρτήσεις των οποίων οι παράγωγοι να είναι είτε εξ ολοκλήρου αρνητικές είτε εξ ολοκλήρου θετικές (είτε, σαφώς, να μη διατηρούν πρόσημο). Ας πούμε, ένας τρόπος να συμπληρώσουμε τον παραπάνω πίνακα είναι αυτός που φαίνεται στον Πίνακα \ref{tab:2}
	
	\begin{table}[!htb]
		\centering
		\def\arraystretch{1.2}
		\begin{tabular}{c|c|c}
			\diagbox{$f$}{$f'$} & $+$ & $-$ \\\hline
			$+$ & $e^x$ & $-e^{-x}$ \\\hline
			$-$ & $e^{-x}$ & $-e^x$ 
		\end{tabular}
		\caption{Ευκολάκι\ldots}
		\label{tab:2}
	\end{table}

	Κι εδώ γεννιέται ένα γενικότερο ερώτημα:
	
	\begin{quest}\label{quest:sign}
		Αν $(s_0,s_1,s_2,\ldots)\in\{-1,1\}^\N$ είναι μία άπειρη ακολουθία από $-1$ και 1, υπάρχει άπειρες φορές παραγωγίσιμη συνάρτηση $f:\R\to\R$ έτσι ώστε $\sgn\left(f^{(n)}\right)s_n=1$, όπου $\sgn(f)=1$ αν $f$ θετική και $\sgn(f)=-1$ αν $f$ αρνητική σε όλο το $\R$; Με άλλα λόγια, υπάρχει άπειρες φορές παραγωγίσιμη συνάρτηση $f$ με το πρόσημο των παραγώγων της να δίνεται από την $s_n$;
	\end{quest}

	Ωραίες απορίες έχουμε\ldots
	
	Λοιπόν, επειδή καμία απορία δε λύθηκε με γκρίνια, ας πιάσουμε χαρτί και μολύβι. Μπορούμε, ας πούμε, να βρούμε μία συνάρτηση με παραγώγους που να έχουν πρόσημο όπως η παρακάτω ακολουθία:
	\[s=(1,-1,1,-1,-1,\ldots);\]
	Με ελάχιστη σκέψη μπορούμε να απαντήσουμε καταφατικά: $f(x)=e^{-x}$. Πράγματι, κάθε φορά που παραγωγίζουμε, «κατεβαίνει» ένα μείον και το πρόσημο της συνάρτησής μας αλλάζει. Έτσι, πετυχαίνουμε αυτήν την εναλλάσσουσα συμπεριφορά στο πρόσημο. Το ίδιο μπορούμε να κάνουμε και με την «εξαδέλφη» της παραπάνω:
	\[t=(-1,1,-1,1,-1,\ldots),\]
	αρκεί να επιλέξουμε την $f(x)=-e^{-x}$. Επίσης, αρκετά εύκολα μπορούμε να βρούμε και συναρτήσεις που να έχουν πάντα θετικές ή πάντα αρνητικές παραγώγους (δεν είναι δύσκολο να τις σκεφτείτε). Άρα, μ' αυτά και μ' εκείνα, έχουμε ήδη ξεμπερδέψει με τέσσερις ακολουθίες και μας μένουν άλλες\ldots\  άπειρες. Και μάλιστα, τόσες όσοι και οι πραγματικοί αριθμοί\footnote{Αλήθεια, πού το ξέρουμε αυτό;}. Προφανώς, δε θα βγάλουμε άκρη αν συνεχίσουμε να παίζουμε έτσι με διάφορες ακολουθίες που κατεβάζει το κεφάλι μας, ωστόσο\ldots\ θα συνεχίσουμε να το κάνουμε για λίγο ακόμα. Κι αυτό γιατί, τρώγοντας έρχεται η όρεξη και, εν προκειμένω, παίζοντας έρχετ' η απόδειξη, που λέμε εμείς στα μαθηματικά (;). Διότι, καθώς παίζουμε με διάφορα παραδείγματα, σιγά-σιγά θα αρχίσουμε να αποκτούμε μία διαίσθηση για το πρόβλημά μας η οποία μπορεί να μας οδηγήσει σε μία απόδειξη της εικασίας μας (ότι υπάρχει πάντα τέτοια συνάρτηση) ή σε ένα αντιπαράδειγμα.
	
	Λοιπόν, ας δοκιμάσουμε να βρούμε μία συνάρτηση που οι παράγωγοί της να έχουν το εξής πρόσημο\footnote{Τόση ώρα λέμε «οι παράγωγοι» αλλά συμπεριλαμβάνουμε, ως \textit{μηδενική παράγωγο}, και την ίδια τη συνάρτηση στη λίστα μας.}:
	\[s=(1,1,1,-1,-1,-1,-1\ldots).\]
	Δηλαδή, για να είμαστε και σαφείς, θέλουμε οι $f,f',f''$ να είναι θετικές και όλες οι επόμενες παράγωγοι αρνητικές. Δηλαδή, μιας και εδώ έχουμε και γεωμετρική ερμηνεία, θέλουμε μία συνάρτηση θετική, γνησίως αύξουσα και κυρτή, που κάθε επόμενη παράγωγός της να είναι αρνητική. Αυτό σημαίνει ότι θέλουμε η δεύτερη παράγωγός της να είναι γνησίως φθίνουσα και, ως εκ τούτου, η πρώτη παράγωγός της να είναι κοίλη.
	
	Χμμμ\ldots
	
	Σίγουρα καμία εκθετική δε μας κάνει. Επίσης, δε μας κάνει και καμία πολυωνυμική, γιατί μηδενίζουν οι παράγωγοί της από ένα σημείο και μετά. Σίγουρα όχι κι οι τριγωνομετρικές, γιατί δεν είναι καν μονότονες. Σίγουρα όχι οι λογαριθμικές γιατί είναι κοίλες για βάση μεγαλύτερη του 1 και γνησίως φθίνουσες για βάση μικρότερη του 1. Σίγουρα έχουμε βρεθεί σε αδιέξοδο. Ας πάμε, λοιπόν, να ψάξουμε μία πιο απλή ακολουθία πρώτα:
	\[s_2=(1,-1,-1,\ldots).\]
	Πήραμε, δηλαδή, την $s$ και της κόψαμε τους δύο πρώτους όρους (εξ ου και το όνομα $s_2$). Επομένως, αναζητούμε θετική, φθίνουσα και κοίλη συνάρτηση, με όλες τις επόμενες παραγώγους αρνητικές. Αρχικά, ως θετική και φθίνουσα θα έχει όριο στο $+\infty$ καθώς το σύνολο τιμών της, $f(\R)$, είναι μη κενό και κάτω φραγμένο, κι άρα από το αξίωμα της πληρότητας υπάρχει το $\inf f(\R)$ το οποίο είναι εύκολο (και διαισθητικά προφανές) να αποδείξουμε ότι είναι το όριό της. Βασικά, αντί να το αποδείξουμε, θα παρατηρήσουμε απλώς το Σχήμα \ref{fig:2} που αποτυπώνει ακριβώς αυτήν την ιδέα.
	
	\begin{figure}[!htb]
		\centering
		\begin{tikzpicture}
		\draw[thick, ->] (-1.5,0) -- (6,0);
		\draw[thick, ->] (0,-1) -- (0,4);
		\draw[thick, domain={-1}:{6}, samples=100] plot (\x,{.5+2^(-\x)});
		\draw[dashed] (-1.5,.5) -- (6,.5)node[pos=0,left]{$\inf f(\R)$};
		\end{tikzpicture}
		\caption{Καλό είναι να θέτουμε όρια πού και πού\ldots}
		\label{fig:2}
	\end{figure}

	Έχοντας αυτό κατά νου, μπορούμε να αναρρωτηθούμε το εξής: Μπορεί μια τέτοια συνάρτηση να είναι κοίλη; Ας πάρουμε λίγο το φανταστικό μολύβι μας κι ας προσπαθήσουμε να σχεδιάσουμε μία συνάρτηση $f$ η οποία να ικανοποιεί τα εξής:
	\begin{itemize}
		\item να είναι θετική,
		\item να είναι γνησίως φθίνουσα,
		\item να είναι κοίλη.
	\end{itemize}
	Αφού είναι θετική και γνησίως φθίνουσα, όπως είπαμε, θα πρέπει να έχει όριο στο $+\infty$. Τώρα, αφού είναι κοίλη, η πρώτη της παράγωγος είναι κι αυτή φθίνουσα, άρα, δεδομένου ότι κι η $f$ είναι φθίνουσα, ο ρυθμός μεταβολής της $f$ γίνεται ολοένα και πιο αρνητικός. Προσέξτε, εδώ αναφερόμαστε στην ουσία στον ρυθμό μεταβολής \textit{του ρυθμού μεταβολής}. Δηλαδή, η $f$ καθώς προχωράμε δεξιά, όχι μόνο κατεβαίνει, αλλά κατεβαίνει όλο και πιο απότομα. Συνεπώς, αν υποθέσουμε για ευκολία ότι $\lim\limits_{x\to+\infty}f(x)=0$, τότε, αν $f(1)=10$ και $f(101)=5$, θα πρέπει να ισχύει ότι $f(201)<0$, γιατί η $f$ θα έχει κατέβει από το $101$ στο $201$ (που απέχουν 100 μονάδες) περισσότερο από ό,τι είχε κατέβει από το $1$ έως το $101$ (που απέχουν επίσης 100 μονάδες).
	
	Πιο αυστηρά, γνωρίζουμε από το σχολείο ότι μία παραγωγίσιμη συνάρτηση είναι κοίλη ακριβώς όταν η παράγωγός της είναι γνησίως φθίνουσα. Ωραία, ας πάρουμε τώρα τον αριθμό $f(1)$. Αυτός είναι ένας ωραιότατος θετικός αριθμός. Επειδή η $f$ είναι γνησίως φθίνουσα, συνεχής και $\lim\limits_{x\to+\infty}f(x)=0$, έπεται ότι $f([1,+\infty))=(0,f(1)]$. Δηλαδή, η $f$ θα παίρνει μετά το 1 όλες τις θετικές τιμές που είναι μικρότερες του $f(1)$ (και, μάλιστα, ακριβώς μία φορά).
	
	Συνεπώς, υπάρχει (μοναδικό) $x_0>1$ έτσι ώστε:
	\[f(x_0)=\frac{f(1)}{2}.\]
	Τώρα, θεωρούμε τον αριθμό $y_0=x_0+x_0-1=2x_0-1$. Με άλλα λόγια, απομακρυνόμαστε προς τα δεξιά από το $x_0$ τόσο όσο απέχει το $x_0$ από το 1 ($x_0-1$, δηλαδή). Τώρα, έχουμε δύο όμορφα κλειστά διαστήματα, το $[1,x_0]$ και το $[x_0,y_0]$ που έχουν ακριβώς το ίδιο μήκος (κρατήστε το αυτό για μετά). Ε, κλειστά διαστήματα έχουμε, παραγωγίσιμη συνάρτηση έχουμε, δεν κάνουμε κανένα ΘΜΤ; (πολύ «σχολικό» αυτό το σκεπτικό, αλλά δουλεύει) Λοιπόν, από εφαρμογή του Θεωρήματος Μέσης Τιμής στα διαστήματα $[1,x_0]$ και $[x_0,y_0]$ βρίσκουμε $\xi_1\in(1,x_0)$ και $\xi_2\in(x_0,y_0)$ έτσι ώστε να ισχύουν τα εξής:
	\begin{align}
	f'(\xi_1)=\frac{f(x_0)-f(1)}{x_0-1},\label{eq:mvt1}\\
	f'(\xi_2)=\frac{f(y_0)-f(x_0)}{y_0-x_0}.\label{eq:mvt2}
	\end{align}
	Τώρα, αφού η $f$ είναι κοίλη, η $f'$ είναι γνησίως φθίνουσα, άρα:
	\begin{equation}\label{eq:mvts}
	\xi_1<\xi_2\Rightarrow f'(\xi_1)<f'(\xi_2)\overset{\eqref{eq:mvt1}}{\underset{\eqref{eq:mvt2}}{\Rightarrow}}\frac{f(x_0)-f(1)}{x_0-1}<\frac{f(y_0)-f(x_0)}{y_0-x_0}.
	\end{equation}
	Κι εδώ κάπου να θυμηθούμε ότι $y_0-x_0=2x_0-1-x_0=x_0-1>0$ και άρα:
	\[\eqref{eq:mvts}\Rightarrow f(x_0)-f(1)<f(y_0)-f(x_0)\Rightarrow f(y_0)<2f(x_0)-f(1).\]
	Κι εδώ θα θυμηθούμε και κάτι ακόμα:
	\[f(x_0)=\frac{f(1)}{2}\Rightarrow2f(x_0)=f(1)\Rightarrow2f(x_0)-f(1)=0.\]
	Άρα $f(y_0)<0$, άτοπο, καθώς η $f$ έχει υποτεθεί θετική!
	
	Συνεπώς, δεν υπάρχει θετική, γνησίως φθίνουσα και κοίλη συνάρτηση, άρα ούτε και συνάρτηση που να έχει παραγώγους με πρόσημα όπως η ακολουθία $s_2$, άρα ούτε και συνάρτηση με παραγώγους με πρόσημα όπως και $s$ (αλλιώς η δεύτερη παράγωγός της θα είχε πρόσημα όπως η $s_2$). Άρα, η απάντηση στην Απορία \ref{quest:sign} είναι ηχηρά αρνητική! Κρίμα, και φαινόταν εύλογη γενίκευση\ldots
	
	\section{Λιγότερο γενικά;}
	Εντάξει, ίσως παραγενικεύσαμε στα παραπάνω. Ωστόσο, όταν μία γενίκευση αποτυγχάνει, τι πιο φυσιολογικό από το να αναρρωτηθούμε το εξής:
\end{document}